\SourcePage{34}

按 2.6.2(i)，$\mathcal{D}_X$-代数态射 $h$ 诱导 $\mathbb{C}$-代数态射
\[
  h_X:\mathfrak{z}(X)\longrightarrow\Gamma(\operatorname{Bun}_G,D')
\]

我们将证明 $(\mathfrak{z}(X),h_X)$ 是 2.2.5 中量子化问题的一个解。
在此之前，须先定义 $\mathfrak{z}(X)$ 上的滤链，以及同构
$\sigma_{\mathfrak{z}(X)}:\operatorname{gr}\mathfrak{z}(X)\xrightarrow{\sim}\mathfrak{z}^{cl}(X)$。

$\mathfrak{z}_{\mathfrak{g}}(\mathcal{O})$ 上的典范滤链是
$\operatorname{Aut}\mathcal{O}$-不变的。同构
$\sigma_{\mathfrak{z}(\mathcal{O})}:\operatorname{gr}\mathfrak{z}_{\mathfrak{g}}(\mathcal{O})
\xrightarrow{\sim}\mathfrak{z}^{cl}_{\mathfrak{g}}(\mathcal{O})$（见 2.5.2(ii)）
与 $\operatorname{Aut}\mathcal{O}$-作用相容。因此 $\mathfrak{z}$ 带有水平滤链，
并且有 $\mathcal{D}_X$-代数的同构
\[
  \sigma_{\mathfrak{z}}:\operatorname{gr}\mathfrak{z}\xrightarrow{\sim}\mathfrak{z}^{cl}
\]
它在每条纤维上化为 2.5.2(ii) 的同构 $\sigma_{\mathfrak{z}_x}$。

通过 $\theta_{\mathfrak{z}}:\mathfrak{z}\to H_\nabla(X,\mathfrak{z})_X=\mathfrak{z}(X)_X$，
该滤链的像给出 $\mathfrak{z}(X)_X$ 上的水平滤链，
这等同于 $\mathfrak{z}(X)$ 上的滤链。考虑分次 $\mathcal{D}_X$-代数的满态射
$(\operatorname{gr}\theta_{\mathfrak{z}})\sigma_{\mathfrak{z}}^{-1}:\mathfrak{z}^{cl}\to
\operatorname{gr}\mathfrak{z}(X)_X$。由伴随（见 (22)），
它给出分次 $\mathbb{C}$-代数的满态射
\[
  j:\mathfrak{z}^{cl}(X)=H_\nabla(X,\mathfrak{z}^{cl})\longrightarrow\operatorname{gr}\mathfrak{z}(X).
\]

注意 $h_X$ 与滤链相容，且有交换图
\begin{equation}\tag{28}
\begin{tikzcd}[column sep=3.4em, row sep=2.4em]
\mathfrak{z}^{cl}(X) \arrow[r, hook, "h^{cl}_X"] \arrow[d, "j"]
  & \Gamma(\operatorname{Bun}_G,P) \arrow[d, "\sigma_{\operatorname{Bun}_G}"] \\
\operatorname{gr}\mathfrak{z}(X) \arrow[r, hook, "\operatorname{gr}h_X"]
  & \operatorname{gr}\Gamma(\operatorname{Bun}_G,D')
\end{tikzcd}
\end{equation}

因此 $j$ 是同构，$\operatorname{gr}h_X$（从而 $h_X$）是单射。
定义 $\sigma_{\mathfrak{z}(X)}:\operatorname{gr}\mathfrak{z}(X)\xrightarrow{\sim}
\mathfrak{z}^{cl}(X)$ 为 $\sigma_{\mathfrak{z}(X)}:=j^{-1}$。
三元组 $(\mathfrak{z}(X),h_X,\sigma_{\mathfrak{z}(X)})$ 是 2.2.5 中量子化问题的一个解。

\bdp{2.7.5} 让我们证明定理 2.5.5，并把 2.5.5 中的 $\mathfrak{z}^{(x)}(X)$
与 $\mathfrak{z}(X)$ 作比较。显然 $h_x=h_X\cdot\theta_{\mathfrak{z}_x}$，
其中 $\theta_{\mathfrak{z}_x}:\mathfrak{z}_x\to\mathfrak{z}(X)$ 由命题 2.6.2(ii) 定义。
$\theta_{\mathfrak{z}_x}$ 是满射（见 2.6.2(ii)），且与滤链严格相容
（见 2.7.4 中 $\mathfrak{z}(X)$ 上滤链的定义）。$h_X$

\SourcePage{35}

是单射，且与滤链严格相容（见 2.7.4 末尾）。
于是 $h_x$ 与滤链严格相容（这正是定理 2.5.5），
并且 $h_X$ 诱导滤代数 $\mathfrak{z}(X)$ 与
$\mathfrak{z}^{(x)}(X):=h_x(\mathfrak{z}_x)$ 之间的同构。

\bdsec{2.8. 水平性} 本小节我们在若干自然的模概形上引入 $\mathcal{D}_X$-结构，
并（模去 4.4.14 中解释的某些细节）证明水平性定理 2.7.3。
读者眼下可以跳过本小节。

在 2.8.1--2.8.2 中我们给出定理 2.7.3 的证明概要。2.8.2 的方法在 2.8.3 中略作改动。
在 2.8.4--2.8.5 中我们解释一些细节，其余细节引向 4.4.14。
在 2.8.6 中我们极简地讨论分歧情形。

\bdp{2.8.1} 让我们构造定理 2.7.3 中的态射 $h$。

回忆 2.5.3--2.5.4 中 $h_x$ 的构造涉及概形 $\operatorname{Bun}_{G,x}$，
即在 $x$ 的形式邻域上平凡化的 $X$ 上 $G$-丛的模概形。
它还涉及 Harish-Chandra 对 $(\mathfrak{g}\otimes K_x,G(\mathcal{O}_x))$
在 $\operatorname{Bun}_{G,x}$ 上的作用，以及该作用提升为
$(\widetilde{\mathfrak{g}\otimes K_x},G(\mathcal{O}_x))$ 在线丛
$\pi_x^*\omega_{\operatorname{Bun}_G}$ 上的作用，
其中 $\pi_x$ 是自然态射 $\operatorname{Bun}_{G,x}\to\operatorname{Bun}_G$。
这些作用来自群归纳概形 $G(K_x)$ 在 $\operatorname{Bun}_{G,x}$ 上的作用，
以及该作用提升为某个中心扩张\footnote{这一扩张在 2.5.3 的注中被提到（而非被定义）。
这对我们所给的概要已经足够。} $\widehat{G(K_x)}$ 在
$\pi_x^*\omega_{\operatorname{Bun}_G}$ 上的作用。

为了构造 $h$，须把上述依赖于 $x$ 的对象组织成族。
用显然的方式定义 $X$ 上的概形 $M$，其纤维在 $x$ 上等于 $\operatorname{Bun}_{G,x}$。
定义 $X$ 上的群概形 $J(G)$ 与 $X$ 上的群归纳概形 $J^{mer}(G)$，
它们在 $x$ 上的纤维分别是 $G(\mathcal{O}_x)$ 与 $G(K_x)$。
$J(G)$ 是函数 $X\to G$ 的射流概形，$J^{mer}(G)$ 是「亚纯射流」的归纳概形。
$J^{mer}(G)$ 作用于 $M$。最后，

\SourcePage{36}

定义一个中心扩张 $\widehat{J}^{mer}(G)$ 及它在
$\pi^*\omega_{\operatorname{Bun}_G}$ 上的作用，
其中 $\pi$ 是自然态射 $M\to\operatorname{Bun}_G$。
这些数据定义好之后，$h:\mathfrak{z}\to\Gamma(\operatorname{Bun}_G,D')_X$
的构造与 $h_x$ 的构造（见 2.5.3--2.5.4）颇为相似。

\bdp{2.8.2} 关键的观察是：\emph{在 $J(G)$、$J^{mer}(G)$、$\widehat{J}^{mer}(G)$、
$M$ 与 $\pi^*\omega_{\operatorname{Bun}_G}$ 上都存在沿 $X$ 的典范联络，
使得 $J^{mer}(G)$ 在 $M$ 上的作用与 $\widehat{J}^{mer}(G)$ 在
$\pi^*\omega_{\operatorname{Bun}_G}$ 上的作用都是水平的。}
由此可得 $h$ 的水平性。

对 $X$-概形 $Y$，记 $\mathcal{J}Y$ 为截面 $X\to Y$ 的射流概形。
众所周知（2.6.3 中有粗略解释），$\mathcal{J}Y$ 带有沿 $X$ 的典范联络
（即 $\mathcal{J}Y$ 是 2.6.1 意义下的 $\mathcal{D}_X$-概形）。
特别地，这适用于 $J(G)=\mathcal{J}(G\times X)$。
若 $\mathcal{F}$ 是 $X$ 上的主 $G$-丛，则 $\pi:M\to\operatorname{Bun}_G$
在 $\mathcal{F}$ 上的纤维等于 $\mathcal{J}\mathcal{F}$，故它带有沿 $X$ 的联络。
可以证明，这些联络来自 $M$ 上沿 $X$ 的联络。

为定义 $M$ 上的联络以及其它联络，使用 Grothendieck 的进路 [Gr68] 较为方便。
按 [Gr68]，$X$-概形 $Z$ 上沿 $X$ 的联络（=可积联络=「分层」）
是一族双射 $\varphi_{\alpha\beta}:\operatorname{Mor}_\alpha(S,Z)\xrightarrow{\sim}
\operatorname{Mor}_\beta(S,Z)$，对每个概形 $S$ 与每对无限接近的「点」
$\alpha,\beta:S\to X$（这里 $\operatorname{Mor}_\alpha(S,Z)$ 是 $\alpha$ 在
$\operatorname{Mor}(S,Z)$ 中的原像，「无限接近」指 $\alpha$ 与 $\beta$ 在
$S_{red}$ 上的限制相同）；双射 $\varphi_{\alpha\beta}$ 要求关于 $S$ 具有函子性，
并满足方程 $\varphi_{\beta\gamma}\varphi_{\alpha\beta}=\varphi_{\alpha\gamma}$。

例如，若 $Z$ 是 $X$ 上概形 $Y$ 的射流概形，则
$\operatorname{Mor}_\alpha(S,Z):=\operatorname{Mor}_X(S'_\alpha,Y)$，
其中 $S'_\alpha$ 是图像 $\Gamma_\alpha\subset S\times X$ 的形式邻域，
态射 $S'_\alpha\to X$ 由投影 $\operatorname{pr}_X:S\times X\to X$ 诱导。
容易证明，若 $\alpha$ 与 $\beta$ 无限接近，则 $S'_\alpha=S'_\beta$，
于是我们得到 $Z$ 上沿 $X$ 的联络。可以证明它与 2.6.3 中定义的联络一致。

\SourcePage{37}

$J^{mer}(G)$、$\widehat{J}^{mer}(G)$ 与 $M$ 上沿 $X$ 的联络用类似方式定义。
$J^{mer}(G)$ 在 $M$ 上的作用与 $\widehat{J}^{mer}(G)$ 在
$\pi^*\omega_{\operatorname{Bun}_G}$ 上的作用的水平性，由定义即可推出。

\bdp{2.8.3} 2.8.2 中描述的方法可作如下改动。回忆 $\mathcal{O}:=\mathbb{C}[[t]]$，
$K:=\mathbb{C}((t))$；$\operatorname{Aut}\mathcal{O}$ 与 $X^\wedge$ 在 2.6.5 中定义。
置 $M^\wedge=M\times_X X^\wedge$。于是 $M^\wedge$ 是四元组
$(x,t_x,\mathcal{F},\gamma_x)$ 的模空间，其中 $x\in X$，$t_x$ 是 $x$ 处的形式参数，
$\mathcal{F}$ 是 $X$ 上的 $G$-torsor，$\gamma_x$ 是 $\mathcal{F}$ 在 $x$
的形式邻域上的截面。群归纳概形 $G(K)$ 作用在 $M^\wedge$ 于任意
$\hat{x}\in X^\wedge$ 上的纤维（事实上该纤维等于 $\operatorname{Bun}_{G,x}$，
其中 $x$ 是 $\hat{x}$ 在 $X$ 中的像，故 $G(K_x)$ 作用在纤维上；
另一方面，与 $\hat{x}$ 对应的 $x$ 处形式参数给出同构 $K_x\xrightarrow{\sim}K$）。
实际上 $G(K)$ 作用在 $M^\wedge$ 上（见 2.8.4），且中心扩张 $\widehat{G(K)}$
作用在 $\hat{\pi}^*\omega_{\operatorname{Bun}_G}$ 上，
其中 $\hat{\pi}$ 是自然态射 $M^\wedge\to\operatorname{Bun}_G$。
该作用诱导态射
\[
  \hat{h}:\mathfrak{z}_{\mathfrak{g}}(\mathcal{O})\longrightarrow
  \Gamma(X^\wedge,\mathcal{O}_{X^\wedge})\otimes\Gamma(\operatorname{Bun}_G,D')
\]
（$\mathfrak{z}_{\mathfrak{g}}(\mathcal{O})$ 的定义见 2.7.2）。

另一方面，2.6.5 中 $\operatorname{Aut}\mathcal{O}$ 在 $X^\wedge$ 上的作用
典范地提升为其在 $M^\wedge$ 上（见 2.8.4）与在层
$\hat{\pi}^*\omega_{\operatorname{Bun}_G}$ 上的作用。
$\operatorname{Aut}\mathcal{O}$ 与 $\widehat{G(K)}$ 在
$\hat{\pi}^*\omega_{\operatorname{Bun}_G}$ 上的作用以显然的方式相容。
因此 $\hat{h}$ 是 $\operatorname{Aut}\mathcal{O}$-等变的。
于是 $\hat{h}$ 诱导水平态射
$h:\mathfrak{z}=\mathfrak{z}_{\mathfrak{g}}(\mathcal{O})_X\to
\Gamma(\operatorname{Bun}_G,D')_X$。

\bdp{2.8.4} 为把 2.8.3 的概要变成定理 2.7.3 的证明，
我们首先给出半直积 $\operatorname{Aut}\mathcal{O}\ltimes G(K)$ 在 $M^\wedge$
上作用的精确定义。设 $R$ 是 $\mathbb{C}$-代数。按定义，$M^\wedge$ 的 $R$-点
是一个三元组 $(\alpha,\mathcal{F},\gamma)$，其中
$\alpha:\operatorname{Spec} R\widehat{\otimes}\mathcal{O}\to X\otimes R$
是 $R$-态射，其微分在 $\operatorname{Spec} R$ 的每一点处都不为零；
$\mathcal{F}$ 是 $X\otimes R$ 上的 $G$-torsor；$\gamma$ 是 $\alpha^*\mathcal{F}$ 的截面。
记 $\Gamma_\alpha$ 为复合
$\operatorname{Spec} R\to\operatorname{Spec} R\widehat{\otimes}\mathcal{O}
\xrightarrow{\alpha}X\otimes R$ 的图像，$\alpha'$ 为 $\alpha$ 诱导的态射
$\operatorname{Spec} R\widehat{\otimes}K\to(X\otimes R)\setminus\Gamma_\alpha$。按

\SourcePage{38}

Beauville 与 Laszlo\footnote{$\Gamma_\alpha\subset X\otimes R$ 的法丛平凡，
故按 2.3.7 可以应用 [BLa95] 的主定理，而不必用其整体化版本。}
（见 2.3.7 与 2.3.4），$M^\wedge$ 的 $R$-点与三元组 $(\alpha,\mathcal{F}',\gamma')$
一一对应，其中 $\alpha$ 如上，$\mathcal{F}'$ 是 $(X\otimes R)\setminus\Gamma_\alpha$
上的 $G$-torsor，$\gamma'$ 是 $\alpha'^*\mathcal{F}'$ 的截面
（当然，$\mathcal{F}'$ 是 $\mathcal{F}$ 的限制，$\gamma'$ 是 $\gamma$ 的限制）。
这一解释表明 $G(R\widehat{\otimes}K)$ 与
$\operatorname{Aut}(R\widehat{\otimes}\mathcal{O})$ 作用在 $M^\wedge(R)$ 上：
$G(R\widehat{\otimes}K)$ 的作用改变 $\gamma'$，
$\operatorname{Aut}(R\widehat{\otimes}\mathcal{O})$ 的作用改变 $\alpha$
（若把 $\alpha$ 换成 $\alpha\circ\varphi$，
$\varphi\in\operatorname{Aut}\operatorname{Spec} R\widehat{\otimes}\mathcal{O}$，
则 $\Gamma_\alpha$ 作为 $X\otimes R$ 的子概形改变，但作为子集不变，
故 $(X\otimes R)\setminus\Gamma_\alpha$ 不变）。
这样我们得到 2.8.3 中提到的 $\operatorname{Aut}\mathcal{O}\ltimes G(K)$
在 $M^\wedge$ 上的作用。

\bdp{2.8.5} 按 2.8.4，$\operatorname{Aut}\mathcal{O}$ 作用在作为
$\operatorname{Bun}_G$ 上概形的 $M^\wedge$ 上。于是 $\operatorname{Aut}\mathcal{O}$
作用在 $\hat{\pi}^*\omega_{\operatorname{Bun}_G}$ 上。
在 2.5.3 中我们提到 $\widehat{G(K_x)}$ 在 $\omega_{\operatorname{Bun}_G}$
到 $\operatorname{Bun}_{G,x}$ 的拉回上的典范作用。
于是 $\widehat{G(K)}$ 作用在 $\hat{\pi}^*\omega_{\operatorname{Bun}_G}$
于任意 $\hat{x}\in X^\wedge$ 上纤维的限制上。如 2.8.3 中所解释，
为完成 2.7.3 的证明，只需证明：
\begin{enumerate}[label=(\roman*),leftmargin=2.6em]
\item 对应于不同 $\hat{x}\in X^\wedge$ 的 $\widehat{G(K)}$ 作用，
  来自 $\widehat{G(K)}$ 在 $\hat{\pi}^*\omega_{\operatorname{Bun}_G}$ 上的一个
  （显然唯一的）作用；
\item 该作用与 $\operatorname{Aut}\mathcal{O}$ 的作用相容。
\end{enumerate}
为证明 (i) 与 (ii)，必须（也几乎充分）定义中心扩张 $\widehat{G(K_x)}$
及其在 $\omega_{\operatorname{Bun}_G}$ 到 $\operatorname{Bun}_{G,x}$
的拉回上的作用。感兴趣的读者可以用例如 [BLa94] 来做。

我们不证明 (i) 与 (ii)，而是在 4.4.14 中对
$\omega_{\operatorname{Bun}_G}$ 的平方根证明类似的陈述
（因为我们需要 $\omega_{\operatorname{Bun}_G}$ 的平方根来表述并证明定理 5.4.5，
那是本工作的主要结果）。更精确地说，对 $\omega_X$ 的任意平方根
$\mathcal{L}$，定义 $\operatorname{Bun}_G$ 上的线丛
$\lambda'_{\mathcal{L}}$，它本质上是 $\omega_{\operatorname{Bun}_G}$ 的平方根
（见 4.4.1）。构造中心扩张\footnote{事实上，该扩张是
$\widehat{G(K_x)}$ 的平方根。} $\widehat{G(K_x)}_{\mathcal{L}}$，
它作用在 $\lambda_{\mathcal{L}}$ 到 $\operatorname{Bun}_{G,x}$ 的拉回上
（见 4.4.7--4.4.8）。从 2.5.4 来的态射
$h_x:\mathfrak{z}_x\to\Gamma(\operatorname{Bun}_G,D')$ 可以

\SourcePage{39}

用该作用自然地定义（见 4.4.12--4.4.13）。
最后，在 4.4.14 中我们对 $\lambda'_{\mathcal{L}}$ 证明上述陈述 (i) 与 (ii)
的类似物，由此推出水平性定理 2.7.3。

\bdp{2.8.6} 设 $\Delta\subset X$ 是有限子概形。记
$\operatorname{Bun}_{G,\Delta}$ 为 $X$ 上在 $\Delta$ 上平凡化的 $G$-丛的叠。
记 $D'$ 为 1.1.6 中的层 $D_{\mathcal{Y},\mathcal{L}^{\lambda}}$，
其中 $\mathcal{Y}=\operatorname{Bun}_{G,\Delta}$，
$\mathcal{L}=\omega_{\operatorname{Bun}_G}$ 的拉回，$\lambda=1/2$。
正如 $\Delta=\emptyset$ 的情形，定义水平态射
\[
  h:\mathfrak{z}_{X\setminus\Delta}\longrightarrow
  \Gamma(\operatorname{Bun}_{G,\Delta},D')\otimes\mathcal{O}_{X\setminus\Delta},
\]
其中 $\mathfrak{z}_{X\setminus\Delta}$ 是 $\mathfrak{z}$ 到 $X\setminus\Delta$ 的限制。
$h$ 诱导单射 $\Gamma(N,\mathcal{O}_N)\to\Gamma(\operatorname{Bun}_{G,\Delta},D')$，
其中 $N=N_\Delta(G)$ 是归纳概形 $N'_\Delta(G)$ 的闭子概形，
而 $N'_\Delta(G)$ 由 $\operatorname{Spec}\mathfrak{z}_{X\setminus\Delta}$
的水平截面构成。

\emph{问题.} 显式描述 $N_\Delta(G)$。

我们将指出该问题的解法中用到的几何对象。
由于不解释解法的细节，读者可以在不知道问题答案的情况下读本小节其余部分，
答案可在 3.8.2 中找到。

对 $n\in\mathbb{Z}_+$，记 $M_{\Delta,n}$ 为三元组的叠，
三元组由一个点 $x\in X$、$X$ 上的 $G$-丛 $\mathcal{F}$、
以及 $\mathcal{F}$ 在 $\Delta+nx$ 上的平凡化组成
（这里我们把 $X$ 的有限子概形与有效除子等同，故 $\Delta+nx$ 有意义）。
$M_{\Delta,n}$ 是代数叠，且
$M_\Delta:=\varprojlim_n M_{\Delta,n}$ 是 $X$ 上的概形。

\emph{注.} 设 $M_{\Delta,x}$ 是 $M_\Delta$ 在 $x\in X$ 上的纤维。
若 $x\in X\setminus\Delta$，则 $M_{\Delta,x}$ 是在 $\Delta$ 与 $x$ 的形式邻域上
平凡化的 $G$-丛的模概形。若 $x\in\Delta$，则
$M_{\Delta,x}=M_{\Delta\setminus\{x\},x}$。

考虑如下定义的「同余子群」概形 $\underline{\mathrm{G}}_\Delta$。
它是平坦于 $X$ 的概形，且对任意平坦于 $X$ 的概形 $S$ 有
\[
  \operatorname{Mor}_X(S,\underline{\mathrm{G}}_\Delta)
  =\{\,f:S\to G \text{ 满足 } f|_{\Delta_S}=1\,\},
\]
其中 $\Delta_S$ 是 $\Delta$ 在 $S$ 中的原像。
$\underline{\mathrm{G}}_\Delta$ 是 $X$ 上的群概形。
$X$ 上在 $\Delta$ 上平凡化的 $G$-丛等同于
$\underline{\mathrm{G}}_\Delta$-丛（这

\SourcePage{40}

在把 $G$-丛与 $\underline{\mathrm{G}}_\Delta$-丛看作平展拓扑下的 torsor 时
就清楚了）。于是 $\operatorname{Bun}_{G,\Delta}$ 是
$\underline{\mathrm{G}}_\Delta$-丛的叠。

可以证明，若 $D\subset X$ 是有限子概形，且 $\Delta+D$ 按除子意义理解，
则对每个平坦于 $X$ 的概形 $S$
\[
  \operatorname{Mor}_X(S,\underline{\mathrm{G}}_{\Delta+D})
  =\{\,f\in\operatorname{Mor}_X(S,\underline{\mathrm{G}}_\Delta)
  \text{ 满足 } f|_{D_S}=1\,\}.
\]

因此，$X$ 上在 $\Delta+D$ 上平凡化的 $G$-丛等同于在 $D$ 上平凡化的
$\underline{\mathrm{G}}_\Delta$-丛。于是 $M_\Delta$ 是三元组的模概形，
三元组由一个点 $x\in X$、$X$ 上的 $\underline{\mathrm{G}}_\Delta$-丛、
以及它在 $x$ 的形式邻域上的平凡化组成。
现在容易定义 $\mathcal{J}^{mer}(\underline{\mathrm{G}}_\Delta)$
在 $M_\Delta$ 上的典范作用，其中
$\mathcal{J}^{mer}(\underline{\mathrm{G}}_\Delta)$ 是截面
$X\to\underline{\mathrm{G}}_\Delta$ 的「亚纯射流」的群归纳概形。
$\mathcal{J}^{mer}(\underline{\mathrm{G}}_\Delta)$ 与 $M_\Delta$ 都带有
沿 $X$ 的联络，且上述作用是水平的。如此等等……

\emph{注}

\begin{enumerate}[label=(\roman*),leftmargin=2.6em]
\item 若 $\Delta\neq\emptyset$，则 2.8.3 的方法不能避免使用 $X$ 上的群归纳概形。
\item 与无穷维概形及归纳概形（如 $M_\Delta$ 或
  $\mathcal{J}^{mer}(\underline{\mathrm{G}}_\Delta)$）相关的陷阱是存在的。
  这里有一例。态射
  $\underline{\mathrm{G}}_\Delta\to\underline{\mathrm{G}}
  :=\underline{\mathrm{G}}_\emptyset=G\times X$ 诱导
  $f:\mathcal{J}^{mer}(\underline{\mathrm{G}}_\Delta)\to
  \mathcal{J}^{mer}(\underline{\mathrm{G}})$。
  该 $f$ 诱导任意点 $x\in X$ 上纤维的同构
  （$\mathcal{J}^{mer}(\underline{\mathrm{G}}_\Delta)$ 在 $x$ 上的纤维是
  $G(K_x)$，它不依赖 $\Delta$）。但若 $\Delta\neq\emptyset$，
  则 $f$ 不是同构，甚至不是单态射。
\end{enumerate}

\bdsec{2.9. \texorpdfstring{$\mathfrak{z}_{\mathfrak{g}}(\mathcal{O})$}{z(O)} 的交换性}
代数 $\mathfrak{z}_{\mathfrak{g}}(\mathcal{O})$ 与
$\mathfrak{z}_x=\mathfrak{z}_{\mathfrak{g}}(\mathcal{O}_x)$ 在 2.5.1 与 2.7.2
中定义（它们当然是同构的）。Feigin 与 Frenkel 在 [FF92] 中证明了
$\mathfrak{z}_{\mathfrak{g}}(\mathcal{O})$ 交换。
本小节我们给出 $\mathfrak{z}_{\mathfrak{g}}(\mathcal{O})$ 交换性的两个证明：
整体的一个（见 2.9.1--2.9.2）与局部的一个（见 2.9.3--2.9.5）。
后者实际上是 [FF92] 中原证明的一个版本。

\SourcePage{41}

读者眼下可以跳过本小节。本文其余部分不会用到 2.9.1--2.9.2。

\bdp{2.9.1} 让我们证明
\begin{equation}\tag{29}
  [h_x(\mathfrak{z}_x),h_y(\mathfrak{z}_y)]=0
\end{equation}
（$h_x:\mathfrak{z}_x\to\Gamma(\operatorname{Bun}_G,D')$ 的定义见 2.5.4）。
由于 $\mathfrak{z}_x$ 是 $\mathcal{O}_X$-代数
$\mathfrak{z}=\mathfrak{z}_{\mathfrak{g}}(\mathcal{O})_X$ 在 $x$ 处的纤维，
而 $h_x$ 来自 $\mathcal{O}_X$-代数态射
$h:\mathfrak{z}\to\mathcal{O}_X\otimes\Gamma(\operatorname{Bun}_G,D')$，
只需对 $x\neq y$ 证明 (29)。记 $\operatorname{Bun}_{G,x,y}$ 为 $X$ 上在 $x$ 与 $y$
的形式邻域上平凡化的 $G$-丛的模概形。$G(K_x)\times G(K_y)$ 作用在
$\operatorname{Bun}_{G,x,y}$ 上。特别地，Harish-Chandra 对
$\big((\mathfrak{g}\otimes K_x)\times(\mathfrak{g}\otimes K_y),
G(\mathcal{O}_x)\times G(\mathcal{O}_y)\big)$ 作用在
$\operatorname{Bun}_{G,x,y}$ 上。该作用典范地提升为
$\big(\widetilde{(\mathfrak{g}\otimes K_x)\times(\mathfrak{g}\otimes K_y)},
G(\mathcal{O}_x)\times G(\mathcal{O}_y)\big)$
在 $\omega_{\operatorname{Bun}_G}$ 到 $\operatorname{Bun}_{G,x,y}$
的拉回上的作用，其中 $\mathbf{1}_x\in\mathfrak{g}\otimes K_x$ 与
$\mathbf{1}_y\in\mathfrak{g}\otimes K_y$ 按乘以 2 作用，
$G(\mathcal{O}_x)\times G(\mathcal{O}_y)$ 按显然方式作用。
$G(\mathcal{O}_x)\times G(\mathcal{O}_y)$ 在 $\operatorname{Bun}_{G,x,y}$
上的作用满足 1.2.6 中的条件 (16)，且商叠等于 $\operatorname{Bun}_G$。
于是按 1.2.5 有典范态射
$h_{x,y}:\mathfrak{z}_x\otimes\mathfrak{z}_y\to
\Gamma(\operatorname{Bun}_G,D')$。
它在 $\mathfrak{z}_x$ 与 $\mathfrak{z}_y$ 上的限制分别等于 $h_x$ 与 $h_y$。
于是 (29) 显然成立。

\bdp{2.9.2} 让我们证明 $\mathfrak{z}_{\mathfrak{g}}(\mathcal{O})$ 的交换性。
设 $a\in[\mathfrak{z}_{\mathfrak{g}}(\mathcal{O}),\mathfrak{z}_{\mathfrak{g}}(\mathcal{O})]$，
$a\neq 0$。若 $x=y$，则 (29) 表明 $h_x(\mathfrak{z}_x)$ 交换。
于是对任意 $X$、$x\in X$ 与同构 $f:\mathcal{O}\xrightarrow{\sim}\mathcal{O}_x$，
有 $h_x(f_*(a))=0$。设
$\bar{a}\in\mathfrak{z}^{cl}_{\mathfrak{g}}(\mathcal{O})$ 是 $a$ 的主象征。
则对任意如上 $X$、$x$、$f$，有 $h_x^{cl}(f_*(\bar{a}))=0$
（$h_x^{cl}:\mathfrak{z}^{cl}_x\to\Gamma(\operatorname{Bun}_G,P)
=\Gamma(T^*\operatorname{Bun}_G,\mathcal{O})$ 的定义与几何描述见 2.4.3）。
这意味着，把 $\bar{a}$ 看作 $\mathfrak{g}^*\otimes\omega_{\mathcal{O}}$
上的多项式函数（见 2.4.1）时，它具有如下性质：对任意如上 $X$、$x$，
任意在 $x$ 的形式邻域上平凡化的 $X$ 上 $G$-丛 $\mathcal{F}$，
以及任意同构 $\mathcal{O}_x\xrightarrow{\sim}\mathcal{O}$，
$\bar{a}$ 在映射
\[
  H^0(X,\mathfrak{g}^*_{\mathcal{F}}\otimes\omega_X)\longrightarrow
  \mathfrak{g}^*\otimes\omega_{\mathcal{O}_x}\xrightarrow{\sim}
  \mathfrak{g}^*\otimes\omega_{\mathcal{O}}
\]
的像上的限制为零。存在 $n$，使得 $\bar{a}$ 来自
$\mathfrak{g}^*\otimes(\omega_{\mathcal{O}}/\mathfrak{m}^n\omega_{\mathcal{O}})$
上的函数，其中 $\mathfrak{m}$ 是 $\mathcal{O}$ 的极大理想。
选取 $X$ 与 $x$，使得映射

\SourcePage{42}

$H^0(X,\omega_X)\to\omega_{\mathcal{O}_x}/\mathfrak{m}^n_x\omega_{\mathcal{O}_x}$
是满射，并令 $\mathcal{F}$ 为平凡丛。则映射
$H^0(X,\mathfrak{g}^*_{\mathcal{F}}\otimes\omega_X)\to
\mathfrak{g}^*\otimes(\omega_{\mathcal{O}_x}/\mathfrak{m}^n_x\omega_{\mathcal{O}_x})$
是满射，从而 $\bar{a}=0$，即矛盾。

\emph{注.} 设 $\operatorname{Bun}_{G,\Delta}$ 是 $X$ 上在有限子概形
$\Delta\subset X$ 上平凡化的 $G$-丛的叠。为从 (29) 推出
$\mathfrak{z}_{\mathfrak{g}}(\mathcal{O})$ 的交换性，可以使用从 $\mathfrak{z}_x$
（$x\notin\Delta$）到 $\operatorname{Bun}_{G,\Delta}$ 上扭微分算子环的自然同态。
于是，不必像上面证明中那样选取 $(X,x)$，可以固定 $(X,x)$ 并取 $\Delta$ 足够大。

\bdp{2.9.3} 记 $\bar{\mathfrak{z}}$ 为完备扭泛包络代数
$\overline{U}'(\mathfrak{g}\otimes K)$ 的中心，
$K:=\mathbb{C}((t))\supset\mathbb{C}[[t]]=\mathcal{O}$。在 [FF92] 中，
Feigin 与 Frenkel 从自然同态
$f:\bar{\mathfrak{z}}\to\mathfrak{z}_{\mathfrak{g}}(\mathcal{O})$
的满射性推出 $\mathfrak{z}_{\mathfrak{g}}(\mathcal{O})$ 的交换性。
我们将给出 $f$ 满射性的一个证明：它可视为 [FF92] 中证明的几何版本，
也可视为 2.4.2 末尾注的「量子化」。
与 [FF92] 及 2.4.2 的关系将在 2.9.7 与 2.9.8 中说明。

\emph{注.} 在中心扩张 $\mathfrak{g}\otimes K$ 的定义中（见 2.5.1），
从而在 $\bar{\mathfrak{z}}$ 与 $\mathfrak{z}_{\mathfrak{g}}(\mathcal{O})$ 的定义中，
我们使用了由 (18) 定义的「临界」双线性形式 $c$。
在 $f$ 满射性的证明中，可以假设 $c$ 是 $\mathfrak{g}$ 上任一不变对称双线性形式，
且 $\mathfrak{g}$ 是任一有限维李代数。另一方面，已知若 $\mathfrak{g}$ 单且 $c$ 非临界，
则相应的代数 $\mathfrak{z}_{\mathfrak{g}}(\mathcal{O})$ 平凡（见 ???）。
% [Tutor Note: 原书 p.42 此处交叉引用印作「???」，底本如此，保留原样。]

\bdp{2.9.4} 我们需要 [BD94] 中
$\overline{U}':=\overline{U}'(\mathfrak{g}\otimes K)$ 的解释。
记 $U'$ 为 $\mathfrak{g}\otimes K$ 的非完备扭泛包络代数。
对 $n\ge 0$，设 $I_n$ 是 $U'$ 中由
$\mathfrak{g}\otimes\mathfrak{m}^n\subset\mathfrak{g}\otimes\mathcal{O}\subset U'$
生成的左理想。按定义，$\overline{U}':=\varprojlim_n U'/I_n$。
设 $U'_k$ 是 $U'$ 的标准滤链，$\overline{U}'_k$ 是 $U'_k$ 在
$\overline{U}'$ 中的闭包，即
$\overline{U}'_k:=\varprojlim_n U'_k/I_{n,k}$，$I_{n,k}:=I_n\cap U'_k$。
[BD94] 的主定理把对偶空间 $(U'_k/I_{n,k})^*$ 与某个

\SourcePage{43}

拓扑向量空间 $\Omega_{n,k}$ 等同。于是
\begin{equation}\tag{30}
  U'_k/I_{n,k}=(\Omega_{n,k})^*,\qquad \overline{U}'_k=(\Omega_k)^*
\end{equation}
其中 $\Omega_k=\varinjlim_n\Omega_{n,k}$，$*$ 表示拓扑对偶。

为定义 $\Omega_{n,k}$，需要一些记号。记 $\mathcal{O}_r$
（相应地，$\omega^{\mathcal{O}}_r$）为 $r$ 个 $\mathcal{O}$
（相应地，$\omega_{\mathcal{O}}$）的完备张量积。置
$\omega^K_r=\omega^{\mathcal{O}}_r\widehat{\otimes}_{\mathcal{O}_r}K_r$，
其中 $K_r$ 是 $\mathcal{O}_r$ 的分式域。我们把 $\mathcal{O}_r$ 与
$\mathbb{C}[[t_1,\ldots,t_r]]$ 等同，并把 $\omega^K_r$ 的元素写成
$f(t_1,\ldots,t_r)\,dt_1\ldots dt_r$，其中 $f$ 属于
$\mathbb{C}[[t_1,\ldots,t_r]]$ 的分式域。

\emph{定义.} $\Omega_{n,k}$ 是 $(k+1)$-元组 $(w_0,\ldots,w_k)$ 的集合，
其中 $w_r\in(\mathfrak{g}^*)^{\otimes r}\otimes\omega^K_r$，且满足
\begin{enumerate}[label=\arabic*),leftmargin=2.4em]
\item $w_r$ 关于对称群 $S_r$ 的作用不变（$S_r$ 同时作用在
  $(\mathfrak{g}^*)^{\otimes r}$ 与 $\omega^K_r$ 上）；
\item $w_r$ 在超平面 $t_i=0$（$1\le i\le r$）上有阶 $\le n$ 的极点，
  在超平面 $t_i=t_j$（$1\le i<j\le r$）上有阶 $\le 2$ 的极点，
  且没有其它极点；
\item 若 $w_r=f_r(t_1,\ldots,t_r)\,dt_1\ldots dt_r$，$r\ge 2$，则
\begin{equation}\tag{31}
\begin{aligned}
  f_r(t_1\ldots,t_r)&=\frac{f_{r-2}(t_1,\ldots,t_{r-2})\otimes c}
    {(t_{r-1}-t_r)^2}\\
  &\quad+\frac{\varphi^*(f_{r-1}(t_1,\ldots,t_{r-1}))}{t_{r-1}-t_r}+\cdots
\end{aligned}
\end{equation}
\end{enumerate}
这里 $c\in\mathfrak{g}^*\otimes\mathfrak{g}^*$ 是中心扩张
$\mathfrak{g}\otimes K$ 的定义中用到的双线性形式，
$\varphi^*:(\mathfrak{g}^*)^{\otimes(r-1)}\to(\mathfrak{g}^*)^{\otimes r}$
是映射 $\varphi:\mathfrak{g}^{\otimes r}\to\mathfrak{g}^{\otimes(r-1)}$ 的对偶，
后者由
$\varphi(a_1\otimes\ldots\otimes a_r)
=a_1\otimes\ldots\otimes a_{r-2}\otimes[a_{r-1},a_r]$ 给出；
而 (31) 中的点号表示在超平面 $t_{r-1}=t_r$ 的一般点处没有极点的表达式。

$\Omega_{n,k}$ 上的拓扑由嵌入
\[
  \Omega_{n,k}\hookrightarrow\prod_{0\le r\le k}
  (\mathfrak{g}^*)^{\otimes r}\otimes\Omega^{\mathcal{O}}_r
\]
诱导，该嵌入由 $(w_0,\ldots,w_k)\mapsto(\eta_0,\ldots,\eta_k)$ 给出，其中
\[
  \eta_r=\prod_i t_i^n\cdot\prod_{i<j}(t_i-t_j)^2\cdot w_r.
\]
